\section{Hypothesis tests and confidence sets: complete proofs}
\label{sec:annex-p15-hypothesis-tests-confidence}

\begin{annexitem}{annex:p15.test.randomization}{thm:p15-test-randomization}{Measurable rejection and its exact probability}
\begin{proof}
The map \((x,u)\mapsto\varphi(x)-u\) is measurable, so
\(\{(x,u):u\leq\varphi(x)\}\) is a measurable product-space
event. Its section at \(x\) is an interval in \((0,1)\)
of Lebesgue length \(\varphi(x)\). Product integration gives
\[
 (P_\theta\times\lambda_{01})\{U\leq\varphi(X)\}
             =\int_{\mathsf X}\int_0^1
                    \mathbf1_{\{u\leq\varphi(x)\}}\,du\,dP_\theta(x)
             =E_\theta\varphi(X).
\]
For value zero the section is empty; for value one it is
all of \((0,1)\). Thus both endpoints have the stated
deterministic behavior.
\end{proof}
\end{annexitem}

\begin{annexitem}{annex:p15.likelihood.threshold}{thm:p15-likelihood-threshold}{All threshold and atom cases}
\begin{proof}
The densities are finite \(\mu\)-a.e. because their
integrals are finite. Altering their infinite-value
null sets preserves both laws. Also \(P_0(f_0=0)=0\);
therefore \(L\) is finite \(P_0\)-a.s.
The upper tail \(H(t)=P_0(L>t)\), \(t\geq0\),
decreases to zero as \(t\) tends to infinity.
For \(0<\alpha<1\), the set
\(\{t\geq0:H(t)\leq\alpha\}\) is thus nonempty
and contains a finite point. Let \(c\) be its infimum.

For every \(m\geq1\), monotonicity and the definition
of the infimum give \(H(c+1/m)\leq\alpha\).
The events \(\{L>c+1/m\}\) increase to \(\{L>c\}\),
so continuity of probability gives \(H(c)\leq\alpha\).
If \(c>0\), take numbers \(t_m<c\) increasing to
\(c\). Their tails exceed \(\alpha\), and the
events decrease to \(\{L\geq c\}\), giving
\(P_0(L\geq c)\geq\alpha\).
If \(c=0\), this last inequality holds because
\(L\geq0\) and \(\alpha<1\).

The difference between the two tails is exactly
\(q=P_0(L=c)\). When \(q>0\), the stated
\(\gamma\) lies in \([0,1]\) and fills the
remaining null probability. When \(q=0\),
the inequalities force \(P_0(L>c)=\alpha\)
and no boundary randomization is needed.
The strict likelihood region on \(\{f_0>0\}\)
is \(\{L>c\}\), and the extra region
\(\{f_0=0<f_1\}\) has null probability zero.
Thus the displayed density test has size
exactly \(\alpha\).
Finally \(E_0\mathbf1_{\{f_0=0\}}=0\)
and \(E_0 1=1\), proving the endpoint
size assertions.
\end{proof}
\end{annexitem}

\begin{annexitem}{annex:p15.neyman.pearson}{thm:p15-neyman-pearson}{The global comparison, all equality conditions and level endpoints}
\begin{proof}
On \(\{f_1>cf_0\}\), the test \(\varphi^*\) is 1,
so \((\varphi^*-\psi)(f_1-cf_0)\geq0\).
On \(\{f_1<cf_0\}\), it is zero, and both factors
are nonpositive. On the equality set their
product is zero. The integrand is therefore
nonnegative everywhere and is bounded in
absolute value by the integrable \(f_1+cf_0\).
Expanding its integral gives
\[
 \int(\varphi^*-\psi)(f_1-cf_0)\,d\mu
             =E_1\varphi^*-E_1\psi-c(\alpha-E_0\psi).
\]
Rearranging gives the stated identity:
the power difference is this nonnegative
integral plus \(c(\alpha-E_0\psi)\).
The latter is also nonnegative because \(c\geq0\) and
\(E_0\psi\leq\alpha\). Hence \(\varphi^*\)
is most powerful.

Equality of powers holds exactly when both
nonnegative terms vanish. A nonnegative
measurable function with integral zero
vanishes a.e.: apply the integral bound
to each set where it exceeds \(1/m\).
On the strict positive region this forces
\(\psi=1\) a.e., and on the strict
negative region it forces \(\psi=0\) a.e.
These two conditions conversely make the
integral zero. The other term vanishes
exactly when \(c(\alpha-E_0\psi)=0\).
This proves both directions of every
equality condition.
Since a most-powerful test must tie the
already optimal \(\varphi^*\), these
conditions characterize all of them.
The boundary may carry a nonconstant
measurable randomization: only the size
constraint and the displayed strict-region
conditions matter. At \(c=0\), power is
already 1 on the support of \(P_1\),
and spending all remaining size on
zero-\(P_1\)-density points is unnecessary.

At level zero, \(\int\psi f_0\,d\mu=0\)
forces \(\psi=0\) \(\mu\)-a.e. on
\(\{f_0>0\}\). Its power is thus at
most \(P_1(f_0=0)\), attained by
\(\varphi_0\). Equality is precisely
\(\psi=1\) \(P_1\)-a.e. on that
remaining set. At level one, every
test is feasible and its power is
at most 1, attained by \(\varphi_1\).
Its power is 1 exactly when
\(\int(1-\psi)\,dP_1=0\), equivalently
\(\psi=1\) \(P_1\)-a.e.
This includes singular laws and common
zero-density regions without dividing
by zero or using an infinite threshold
as an ordinary real multiplier.

For indicator tests of equal size, the
same proof splits the integral into
the portions of the two critical
regions outside their intersection.
Their \(P_0\) masses agree, and the
likelihood inequalities compare their
\(P_1\) masses. Thus it contains
the source's full set-difference
comparison as its nonrandomized case.
\end{proof}
\end{annexitem}

\begin{annexitem}{annex:p15.neyman.pearson.variation}{thm:p15-neyman-pearson-variation}{The complete local calculus argument}
\begin{proof}
For \(d\geq2\), let
\[
 I(s)=\int_B\int_{\ell(z)}^{q(z)+s\eta(z)}h(z,t)\,dt\,dz.
\]
These integrals are finite since
\(|h|\leq f_0+kf_1\) is integrable.
The support of \(\eta\) is a compact
subset of \(B\). On it, \(q-\ell\)
has a strictly positive minimum,
and \(\eta\) is bounded.
Thus small positive and negative
\(s\) give valid boundary variations.
Outside this compact support the
integrand difference is zero.
On it, continuity of \(h\) gives a
uniform bound over the compact range
of the perturbed endpoints, and
\[
 \frac1s\int_{q(z)}^{q(z)+s\eta(z)}h(z,t)\,dt
                  \longrightarrow\eta(z)h(z,q(z)).
\]
Its absolute value is bounded by a
constant times \(|\eta(z)|\), an
integrable compactly supported function.
Dominated convergence proves the
derivative formula, for both signs
of \(s\).

If the derivative is zero for every
such \(\eta\), the continuous function
\(z\mapsto h(z,q(z))\) must vanish.
Otherwise it has one strict sign on
a small ball with closure in \(B\).
Choose on that ball a nonnegative
continuous tent function, positive
inside and zero outside.
Its integral against \(h(z,q(z))\)
has the same strict sign, a contradiction.
Thus \(f_0=kf_1\) on the whole patch.
For a region above the boundary,
subtracting the varying lower-limit
integral changes only the sign.
For \(d=1\), the fundamental theorem
of calculus at the variable endpoint
gives the same derivative and condition
directly. No global existence of a
Lagrange multiplier or smooth optimal
boundary has been assumed.
\end{proof}
\end{annexitem}

\begin{annexitem}{annex:p15.binomial.most.powerful}{thm:p15-binomial-most-powerful}{Every likelihood, coupling and finite calculation}
\begin{proof}
The displayed binomial masses sum to 1
by the binomial theorem. Counting the
\(\binom nk\) words with \(k\) successes
gives the independent-Bernoulli representation,
including the two degenerate endpoint laws.
For \(p_0<p_1<1\), the likelihood ratio is
\[
 \frac{b_{n,p_1}(k)}{b_{n,p_0}(k)}
 =\left(\frac{1-p_1}{1-p_0}\right)^n
  \left(\frac{p_1(1-p_0)}{p_0(1-p_1)}\right)^k.
\]
Its base in the final power is strictly
greater than 1. The threshold rule in
the count is therefore exactly a
likelihood-ratio test, with its chosen
size \(\alpha\), and is most powerful
by the preceding theorem.
The boundary \(j,\gamma\) depends only
on \(p_0,\alpha\), not on \(p_1\).
It is consequently most powerful
against every \(p_1\in(p_0,1)\).
Every fixed test's power is a finite
polynomial in \(p\), so passing the
power comparison to \(p_1=1\)
proves that endpoint as well.

For the composite null, use the same
independent uniforms \(U_1,\ldots,U_n\)
and define
\(K(p)=\sum_l\mathbf1_{\{U_l\leq p\}}\).
This count increases pointwise with
\(p\). The test is nondecreasing in
the count, including the value
\(\gamma\in[0,1]\) at its threshold.
Its rejection probability therefore
increases with \(p\), so its size
over \(p\leq p_0\) is exactly
its rejection probability \(\alpha\)
at \(p_0\). Any other level-\(\alpha\)
test for this composite null is
in particular level \(\alpha\)
at \(p_0\). The simple-alternative
comparison just proved bounds its
power for every \(p_1>p_0\).
This proves the stated uniform
optimality with the full null scope.

For \(n=4,p_0=1/2\),
\(P_0(K=4)=1/16\) and
\(P_0(K=3)=4/16=1/4\).
Consequently
\[
 \gamma=\frac{1/10-1/16}{1/4}=\frac3{20},
 \qquad \frac1{16}+\frac3{20}\frac14=\frac1{10}.
\]
At \(p=3/4\), the two relevant
masses are \(81/256\) and
\(108/256\), so the power is
\[
               \frac{81}{256}
                    +\frac3{20}\frac{108}{256}
                   =\frac{243}{640}.
\]
Subtracting from 1 gives type-II
error \(397/640\).
Dropping the randomized boundary
term gives exactly the stated
nonrandomized size and power.

For the two-sided failure, put
\(a=\varphi(0)\), \(b=\varphi(1)\).
Level \(1/2\) at \(p_0=1/2\)
requires \(a+b\leq1\).
At \(p=3/4\), power is
\[
                  \tfrac14a+\tfrac34b
                    =\tfrac14(a+b)+\tfrac12b\leq\tfrac34.
\]
Equality forces \(b=1,a=0\).
At \(p=1/4\), the reversed
calculation forces \(a=1,b=0\)
for its maximum power \(3/4\).
No test has both forms, proving
that there is no uniformly most-powerful
test for these two alternatives.
\end{proof}
\end{annexitem}

\begin{annexitem}{annex:p15.discrete.p.values}{thm:p15-discrete-p-values}{All tail and randomization probabilities}
\begin{proof}
List the statistic's values increasingly
as \(t_1,\ldots,t_m\), with masses
\(q_j=P_0(T=t_j)\), and put
\(s_j=\sum_{\ell>j}q_\ell\).
Then \(\pi(t_j)=s_j+q_j\).
Among values having positive mass,
the event \(\pi(T)\leq a\) is
either empty or an upper tail
beginning at a first such value
\(t_j\). Its probability is
exactly \(s_j+q_j\leq a\).
Zero-mass values do not change
this calculation. This proves
validity for every \(a\in[0,1]\).
If \(a\) lies strictly between
two attainable positive tail
probabilities, the inequality
is strict.

Given \(T=t_j\) with \(q_j>0\),
\(\pi_{\rm r}\) is uniform on
\((s_j,s_j+q_j)\).
These intervals are disjoint
apart from endpoints and together
cover \((0,1)\) apart from finitely
many endpoints. For every Borel
\(B\subseteq(0,1)\), conditioning
on \(T\) gives
\[
 (P_0\times\lambda_{01})\{\pi_{\rm r}(T,U)\in B\}
  =\sum_{j:q_j>0}q_j
           \frac{\lambda(B\cap(s_j,s_j+q_j))}{q_j}
  =\lambda(B).
\]
Hence the randomized p-value
is exactly uniform. Reversing
the order of the statistic
proves the lower-tail assertions.

For the binomial upper-tail value
calculated at \(p_0\), the map
from \(k\) to that value is
nonincreasing. Thus the rejection
event at any fixed p-value cutoff
is an increasing event of \(k\).
The coupling in the previous proof
bounds its probability under
every \(p\leq p_0\) by that
under \(p_0\), which is at
most the cutoff. This gives
composite-null validity.

For \(n=4,p_0=1/2\), the two
inclusive tails are \(1/16\)
and \(5/16\). At count 3,
the strict tail is \(1/16\)
and the atom is \(1/4\),
giving \(1/16+U/4\).
It is at most \(1/10\)
exactly when \(U\leq3/20\).
At count 4 the randomized
value is \(U/16<1/10\);
at counts below 3 its
strict tail already exceeds
\(1/10\). This proves the
exact equivalence with the
specified randomized test.
\end{proof}
\end{annexitem}

\begin{annexitem}{annex:p15.confidence.test.inversion}{thm:p15-confidence-test-inversion}{The pointwise identity and both measurable directions}
\begin{proof}
Joint measurability of \(\varphi\)
makes the membership graph
\(\{(t,x,u):u>\varphi_t(x)\}\)
measurable. For a fixed \(\theta\),
membership of the true target
has, conditionally on \(x\), a
uniform section of length
\(1-\varphi_{\tau(\theta)}(x)\).
Product integration therefore gives
exactly the identity in the
statement, and its level assumption
gives coverage at least \(1-\alpha\).
If the size is exactly \(\alpha\)
at every relevant full parameter,
the coverage is exactly \(1-\alpha\).

Conversely the complement of a
measurable membership graph gives
a jointly measurable indicator
\(\delta_t(x,u)\).
At the true parameter its
expectation is one minus coverage
and hence at most \(\alpha\).
Define
\[
                      \varphi_t(x)
                          =\int_0^1\delta_t(x,u)\,du.
\]
It is jointly measurable in
\((t,x)\). The measurability
of such a parameter integral
follows first for rectangle
indicators, then for nonnegative
simple functions and their
monotone limits; the monotone-class
extension covers the product
sigma-field. Product integration
preserves all rejection probabilities.
Threshold implementation of this
new randomized family need not
give the original random set
point by point, but it preserves
the stated coverage probabilities.

For a nonrandomized acceptance
family, define \(C(x)=\{t:x\in A_t\}\).
Its defining membership equivalence
holds for every point, not only
almost surely. Conversely
\(A_t=\{x:t\in C(x)\}\)
recovers the family. Taking
probabilities gives both directions
of the acceptance/coverage
correspondence.
Because the displayed condition
is imposed for every \(\theta\),
the proof retains all nuisance
parameter values with the same
target. No prior distribution
on those values has been used.
\end{proof}
\end{annexitem}

\begin{annexitem}{annex:p15.neyman.interval.construction}{thm:p15-neyman-interval-construction}{All seven source propositions and endpoint measurability}
\begin{proof}
For a fixed sample point \(x\),
nonemptiness of \(S(x)\) gives
well-defined extended infimum
\(\ell\) and supremum \(u\)
with \(\ell\leq u\).
They cannot both be \(+\infty\)
or both be \(-\infty\), since
\(S(x)\) contains an actual real
number. Every member of \(S(x)\)
lies between the two endpoints.

If \(\ell=u\), their common
value is finite and nonemptiness
forces \(S(x)=\{\ell\}\).
Suppose \(\ell<u\).
For any finite \(t\) strictly
between them, the definitions
of infimum and supremum give
\(a,b\in S(x)\) with \(a<t<b\).
The order-convexity condition
puts \(t\) in \(S(x)\).
If \(\ell\) is finite, choose
a finite \(b\) with \(\ell<b<u\).
All points of \((\ell,b)\)
have just been shown to belong,
so the endpoint condition
puts \(\ell\) in \(S(x)\).
The same argument puts a finite
\(u\) in \(S(x)\).
Infinite endpoints require no
formal point to be added to
the real parameter set.
Thus in every case
\[
                          x\in A_t
                 \quad\Longleftrightarrow\quad
                          \ell(x)\leq t\leq u(x).
\]
This proves the two directions
of the source's membership
equivalence and the full
interval reconstruction.

For rational \(r\),
\(\{\ell<r\}=\{x:\exists t<r,\ x\in A_t\}\)
and
\(\{u>r\}=\{x:\exists t>r,\ x\in A_t\}\).
The additional hypothesis makes
these events measurable, so
\(\ell,u\) are extended-real
measurable functions. Comparison
with a real parameter then also
gives a measurable membership
graph for the closed interval.
The pointwise equivalence above
gives, for every \(\theta\),
\[
 P_\theta\{\ell(X)\leq\tau(\theta)\leq u(X)\}
                            =P_\theta(A_{\tau(\theta)})=q.
\]
At any attainable target value,
its acceptance region is nonempty
because its probability is \(q>0\).
This includes the source's
nonempty-region consequence.

Conversely let a procedure report
a nonempty closed real interval
with measurable extended endpoints,
and define its acceptance sets
by membership. They are measurable,
and the same event identity
transfers coverage to their
probabilities for every full
parameter value. A nonempty
interval contains an actual
real target, establishing the
source's sample-point condition.
Any point between two included
values is included by the
interval inequalities.
If every value in a nonempty
open interval \((a,b)\) is
included, taking limits in
those inequalities includes
both finite endpoints \(a,b\).
Finally the two rational-union
events are exactly the stated
endpoint cuts and are measurable.
These arguments prove the
source's necessary properties,
their converse, and both
infinite-endpoint and singleton
cases with the measurable
conditions made explicit.
\end{proof}
\end{annexitem}

\begin{annexitem}{annex:p15.binomial.tail.roots}{thm:p15-binomial-tail-roots}{The telescoping derivatives and every count endpoint}
\begin{proof}
Differentiate the finite binomial
sum on \(0<p<1\).
The identities
\(j\binom nj=n\binom{n-1}{j-1}\)
and
\((n-j)\binom nj=n\binom{n-1}{j}\)
make adjacent terms cancel.
For the upper tail only the
first positive term remains:
\[
 \frac{d}{dp}\sum_{j=k}^n\binom njp^j(1-p)^{n-j}
              =n\binom{n-1}{k-1}p^{k-1}(1-p)^{n-k}.
\]
The lower-tail sum leaves only
the negative final term, giving
the second derivative in the
statement. Every displayed
factor is positive in its
specified index range, so
the signs are strict.

For \(1\leq k\leq n\),
the upper tail increases
continuously from 0 at
\(p=0\) to 1 at \(p=1\).
For \(0\leq k<n\), the
lower tail decreases continuously
from 1 to 0. Since
\(0<d<1/2\), the intermediate
value theorem and strict
monotonicity give precisely
one root in \((0,1)\) in
each case.

For \(0<k<n\), at the
lower root \(\ell=\ell_n(k)\),
\[
 F_{n,\ell}(k)
       =1-T_{n,\ell}(k)+b_{n,\ell}(k)
       =1-d+b_{n,\ell}(k)>d.
\]
The lower tail decreases
to its value \(d\) at
\(u_n(k)\), so
\(\ell_n(k)<u_n(k)\).
For \(k=0\), the lower
endpoint is 0 and the
upper root is positive.
For \(k=n\), the upper
endpoint is 1 and the
lower root is less than 1.
Thus all intervals are
nonempty. The endpoint
equations
\((1-u_n(0))^n=d\)
and \(\ell_n(n)^n=d\)
give the stated explicit
formulas.
\end{proof}
\end{annexitem}

\begin{annexitem}{annex:p15.binomial.confidence.coverage}{thm:p15-binomial-confidence-coverage}{The precise strict tests and the finite coverage bound}
\begin{proof}
For \(k\geq1\), strict
monotonicity of the upper
tail gives
\[
                   p<\ell_n(k)
                        \quad\Longleftrightarrow\quad
                   T_{n,p}(k)<d.
\]
For \(k=0\), both conditions
are false on \(p\in[0,1]\):
the lower endpoint is 0
and the upper tail is 1.
Similarly for \(k<n\),
\[
                   p>u_n(k)
                        \quad\Longleftrightarrow\quad
                   F_{n,p}(k)<d;
\]
for \(k=n\), both are
false. These equivalences
hold also at \(p=0,1\)
by the endpoint laws.
Their negations prove
the exact pointwise
inversion statement and
retain equality at both
root endpoints.

At a fixed true \(p\),
the inclusive upper and
lower tails are valid
discrete p-values by the
proved finite-tail result.
Each event with tail
strictly below \(d\)
has probability at most
\(d\). The probability
of their union is at
most \(2d=\alpha\).
This proves coverage at
least \(1-\alpha\).
At \(p=0\), \(K=0\)
almost surely and the
reported interval contains
0. At \(p=1\), \(K=n\)
almost surely and it
contains 1, proving
coverage 1 at both
parameter endpoints.

The \(n=4\) calculations
use \(d=1/20\) and the
explicit endpoint formulas.
For \(n=2,k=1\), the
upper-tail equation is
\(2\ell-\ell^2=1/20\),
whose root in \((0,1)\)
is \(1-\sqrt{19/20}\).
The lower-tail equation
is \(1-u^2=1/20\),
giving \(u=\sqrt{19/20}\).
The other intervals for
\(n=2\) are
\[
 [0,1-\sqrt{1/20}],\qquad
 [\sqrt{1/20},1].
\]
Each of these and the
middle interval contains
\(1/2\). Thus the
coverage probability at
that parameter is exactly
1, establishing strict
discrete conservatism.
\end{proof}
\end{annexitem}

\begin{annexitem}{annex:p15.binomial.beta.quantiles}{thm:p15-binomial-beta-quantiles}{The exact beta-tail identity and inverse formulas}
\begin{proof}
For positive integers \(a,b\),
the beta-gamma identity of
\cref{thm:p15-beta-gamma-convolution}
and the integer gamma integral
give
\[
                   B(a,b)=\frac{(a-1)!(b-1)!}{(a+b-1)!}.
\]
The integer gamma formula
also follows directly by
integration by parts from
\(\int_0^\infty e^{-t}\,dt=1\):
each step multiplies the
integral by the preceding
positive integer, with
vanishing endpoint terms.

Let
\[
 I_p(a,b)=\frac1{B(a,b)}
                   \int_0^p t^{a-1}(1-t)^{b-1}\,dt.
\]
Its derivative for
\(a=k,b=n-k+1\) is
exactly the derivative
of \(T_{n,p}(k)\)
proved above. Both
functions are zero at
\(p=0\), so they agree
on \([0,1]\).
Likewise
\[
                 F_{n,p}(k)=1-I_p(k+1,n-k)
                                           \quad(k<n).
\]
The normalized beta
density is positive on
\((0,1)\), and its
cumulative integral is
continuous, strictly
increasing from 0 to 1.
Every quantile appearing
in the statement is
therefore unique.
Inverting the two
identities gives exactly
the lower and upper
beta-quantile formulas.
All their shapes are
positive in the stated
count ranges. The
separately supplied
count endpoints avoid
any appeal to a
zero-shape beta law.
\end{proof}
\end{annexitem}

\begin{annexitem}{annex:p15.confidence.credibility.boundary}{thm:p15-confidence-credibility-boundary}{Coverage frequencies and the two conditional calculations}
\begin{proof}
Fix \(\theta\) and repeat
the data experiment and
any auxiliary randomization
independently. The coverage
indicators \(Y_1,Y_2,\ldots\)
are independent Bernoulli
variables with success
probability
\(q=(P_\theta\times\lambda_{01})\{\tau(\theta)\in C(X,U)\}\).
Their average has mean
\(q\) and variance
\(q(1-q)/m\leq1/(4m)\).
Under the independent repetition law,
the elementary second-moment bound gives
\[
 P\left(\left|
       \frac1m\sum_{j=1}^mY_j-q\right|>\epsilon\right)
                    \leq\frac1{4m\epsilon^2}\longrightarrow0.
\]
This proves the stated
frequency interpretation.
The same calculation
applies to independent
different experiments
when all have the same
exact coverage \(q\).
It does not apply merely
from numerical coverage
claims to arbitrarily
dependent repetitions.
For fixed observed data,
membership of a fixed
parameter in its fixed
reported interval is a
true or false statement;
the repeated-sampling
calculation does not
assign it a posterior
probability.

Now introduce the separate
joint model with a parameter
drawn uniformly on \((0,1)\)
and conditional count
law \(b_{n,p}\). Its
joint density and mass
at \(p,k\) is
\(\binom nkp^k(1-p)^{n-k}\).
Integrating in \(p\)
gives a positive
probability for every
count, and division by
that integral gives
the conditional density
\[
                  \frac{p^k(1-p)^{n-k}}
                              {B(k+1,n-k+1)}.
\]
It is precisely the
stated beta posterior.
This derivation is
ordinary conditional
probability on a
positive-probability
count event.

For \(n=4,k=0\), the
density is \(5(1-p)^4\).
Writing \(u=1-20^{-1/4}\),
its mass in the confidence
interval \([0,u]\) is
\[
                    \int_0^u5(1-p)^4\,dp
                      =1-(1-u)^5=1-20^{-5/4}.
\]
It differs from \(9/10\),
since \(20^{-5/4}\ne1/10\).
For every count, a
central posterior interval
between its \(1/20\)
and \(19/20\) quantiles
has conditional probability
\(18/20=9/10\), by
continuity of the beta
distribution.
At fixed sampling
parameter \(p=0\),
the observed count is
always zero. The lower
posterior quantile is
strictly positive,
so that interval never
contains 0. Its sampling
coverage is therefore
zero, completing the
converse separation. There are
only \(n+1\) counts, and
every lower posterior endpoint
is strictly positive. Their
minimum \(c\) is positive.
For any \(0<p<c\), every
possible reported posterior
interval excludes \(p\), so
its sampling coverage is
also zero at those interior
parameters.
\end{proof}
\end{annexitem}
